%% parts/p04/universal-reductions-supplement.tex
%% Separately authored additions retained from the historical Part IV draft after
%% the exact inherited source through line 8247 replaced its overlapping core.

\section*{Supplementary bridge for universal computation}
\noindent\textit{Reader note.}
The result below is a separately authored bridge between the exact inherited
program numbering and the partial-recursive characterization in
Section~\ref{sec:p04-partial-recursion}.  It supplements, and does not replace
or alter, the exact inherited chapter.

\begin{hbcorollary}{thm:p04-numbering-partial-recursive}{proved=annex:p04-universal.9}{Range of the local numbering}
For every \(k\geq1\) and partial map \(f:\N^k\rightharpoonup\N\), the
following are equivalent:
\begin{enumerate}
  \item some index \(e\) in the effective numbering fixed in
        Section~\ref{sec:universal-reductions-01} satisfies
        \[
          \varphi_e(\vec x)\simeq f(\vec x)
          \qquad\text{for every }\vec x\in\N^k;
        \]
  \item \(f\) is partial recursive in the closure sense of
        Definition~\ref{def:p04-partial-recursive}.
\end{enumerate}
\end{hbcorollary}

\section{Interpreter, specialization, and fixed points}
\label{sec:p04-local-program-completion}

The inherited universal-computation and recursion arguments above leave the
finite compiler implicit.  The following construction supplies it for the
serialization of Definition~\ref{def:p04-program-serialization}.  In
particular, no translation from an unspecified acceptable numbering is
assumed.  Kleene's enumeration and parameter theorems concern his coding of
systems of equations: see \cite[\S65, Theorems XXII--XXIII, pp.~340--342]{Kleene1952}.
His second recursion theorem is the source of the diagonal construction
\cite[\S66, Theorem XXVII, pp.~352--353]{Kleene1952}.  The results here are
proved for the manuscript's single machine numbering.

\begin{hbdefinition}{def:p04-counter-program}{Finite counter procedure}
A \emph{finite counter procedure} has a finite list of counters with values
in \(\N\), a finite list of numbered instructions, and a designated starting
instruction.  Its instructions are: increment a specified counter and go to
a specified instruction; if a specified counter is zero go to one instruction,
and otherwise decrement it and go to another; and return a specified counter.
All targets name instructions in the list.  The input is put in the first
counter and the others start at zero.  A return supplies the designated
counter's value; an infinite run supplies no value.  A literal initialization
by a natural number abbreviates a zeroing loop followed by that many
increment instructions.  This is a finite proof device whose translation
to the fixed Turing-machine model is established next.
\end{hbdefinition}

\needspace{9\baselineskip}
\begin{hblemma}{thm:p04-counter-compilation}{proved=annex:p04.local.counter-compilation}{Counter compilation}
There is a total computable construction taking any finite counter procedure
and producing a valid index in the fixed machine numbering with exactly its
partial numerical input/output behavior.  The construction remains total
when any finite list of natural-number literals is supplied as part of its
input.  Thus its totality concerns production of a finite program; the
produced program may run forever.
\end{hblemma}

\begin{hbtheorem}{thm:p04-local-universal}{proved=annex:p04.local.universal}{Local universal computation}
There is a partial computable \(U:\N^2\rightharpoonup\N\) satisfying
\[
  U(e,x)\simeq\varphi_e(x)
  \qquad(e,x\in\N).
\]
Its algorithm validates and decodes exactly the serialized machine of
Definition~\ref{def:p04-program-serialization}.  Invalid indices,
rejecting runs, and accepting runs with invalid output give divergence.
This supplies a full proof of the function asserted in
Theorem~\ref{thm:universal-computation}.
\end{hbtheorem}

\begin{hbtheorem}{thm:p04-local-parameter}{proved=annex:p04.local.parameter}{Local parameter theorem}
For every fixed \(m,n\geq1\) there is a total computable map
\(\mathsf{par}_n^m:\N^{m+1}\to\N\), taking only valid program codes
as values, such that
\[
  \varphi_{\mathsf{par}_n^m(e,\bar a)}(\bar x)
  \simeq\varphi_e(\bar a,\bar x)
  \qquad(e\in\N,\ \bar a\in\N^m,\ \bar x\in\N^n).
\]
The input \(e\) need not be valid.  The same finite construction, with
\(m,n\) supplied, computes these maps uniformly over positive arities.
This supplies the compiler verification left open in
Theorem~\ref{thm:parameter}.
\end{hbtheorem}

\begin{hbtheorem}{thm:p04-local-recursion}{proved=annex:p04.local.recursion}{Local recursion theorem}
For every \(k\geq1\) and partial computable
\(g:\N^{k+1}\rightharpoonup\N\), there is a valid index \(e\) such that
\[
  \varphi_e(\bar x)\simeq g(e,\bar x)
  \qquad(\bar x\in\N^k).
\]
In particular, if \(f:\N\to\N\) is total computable, some valid index
\(e\) satisfies \(\varphi_e\simeq\varphi_{f(e)}\), as in
Theorem~\ref{thm:recursion-theorem}.  The equality is equality of partial
functions, including their domains, and does not assert \(e=f(e)\).
\end{hbtheorem}

These constructions use only finite syntactic compilation and simulation.
They neither decide whether a program returns nor select an index from an
extensional description of an arbitrary function.  The recursion result
also makes no least-domain claim: Kleene's first recursion theorem concerns
a different, functional fixed-point construction
\cite[\S66, Theorem XXVI and discussion, pp.~348--354]{Kleene1952}.
